\scanpage{70}

% top, item 1
\textcircled{1}\quad $k_*\xrightarrow{\varepsilon_*}A_*^*$, $\unsure{K_*|k_*|\kappa_*|K_X}=(K_n)_{n\geq0}\quad K_n:A^n_n\mathrel{\unsure{\relbar|\longrightarrow}} k$, satisf. $\left\{\begin{aligned}&\text{Stokes}\\&K_0\varepsilon_0=\mathrm{id}\end{aligned}\right.$

% upper, item 2
\textcircled{2}\quad Sur \deleted{\unsure{espace|espa}} espace topologique $X$, faisceau différentiel \inserted{de $S$-\unsure{modules|modules alg.}} \deleted{\unsure{$E$|$\varepsilon$}} $\Omega^*$ qui soit une \uline{résolution} de $k_X$ (p.ex. complexe de DR sur $X$ variété $C^\infty$, $k=\mathbb{R}=S$)

% upper middle, item 3, qualifying clauses
\textcircled{3}\quad Un sous-ensemble $\frac{1}{2}$ simplicial $\unsure{S(X)'|\mathbb S(X)'}$ du complexe singulier de $X$, tel que $S(X)'\longrightarrow S(X)$ \deleted{\unsure{soit}} induise un isom sur les groupes d'homologie entière (p.ex. $X$ une variété $C^\infty$ et $S(X)'$ ens. formé des simplexes $C^\infty$) et il en \unsure{soit de même|soit} si on \deleted{\unsure{en}} \unsure{remplace} $X$ par ouvert (d'une base) et qu'on \unsure{prend} de $S(X)$ et $S(X)'$ les simplexes de supp. $\subset U$.

% sloping left margin at item 3
% slanted
\margin{donc, si on regarde le faisceau de coch. sing. $S^*_{X,k}$ sur $X$, et de $=$ $S'_{X,k}$, alors $S'_{X,k}$ comme $S^*_{X,k}$ est une résolution de $k$}

% middle, item 4, overwritten transition and insertions
\textcircled{4}\quad $\forall\sigma\in S(X)'$, et \deleted{\unsure{dont on se donne}} \inserted{hom. \deleted{\unsure{de $k$-alg.}}\inserted{de complexes}} $k$-augm.

% middle, pullback on germs
\[
\sigma^*:\varinjlim_{\substack{U\ \text{ouvert de }X\\ U\supset\operatorname{supp}\sigma}}\Omega^*(U)\longrightarrow A_n^*
\quad(n=\dim\sigma)
\]

% middle, compatibility assumption
On suppose que si $\alpha:\Delta_m\longrightarrow\Delta_n$, $\sigma\in S'(X)_n$, \deleted{$\unsure{\sigma}:\lvert\Delta_n\rvert\longrightarrow X$} on a commutativité dans

% center lower, restriction square
\diagram{\begin{tikzcd}[ampersand replacement=\&,row sep=large]
\displaystyle\varinjlim_{U\supset\operatorname{supp}\sigma}\Omega^*(U)\arrow[r]\arrow[d,"\text{restriction}"] \& A_n^*\arrow[d,"\sigma_A^*"]\\
\displaystyle\varinjlim_{V\supset\operatorname{supp}\alpha^*(\sigma)}\Omega^*(V)\arrow[r] \& A_m^*
\end{tikzcd}}

% left of restriction square
\margin{\uline{NB} $\operatorname{supp}\alpha^*(\sigma)\subset\operatorname{supp}\sigma$}

% lower, bracketed NB
[ \uline{NB} On peut dire que $\sigma\longmapsto\varinjlim_{U\supset\operatorname{supp}\sigma}\Omega^*(U)$ forme\inserted{$\Omega^*_{S(X)'}$} un \deleted{\unsure{faisc.}} objet semi-simplicial (en fait, un complexe de $S$-modules $k$-augmentés) sur $S(X)'=\mathcal{X}$,

% lower left, triangle beside bracketed NB
\margin{\diagram{\resizebox{\linewidth}{!}{\begin{tikzcd}[ampersand replacement=\&,column sep=tiny,row sep=tiny]
\Delta_n\arrow[dr,"\sigma"'] \& \& \Delta_m\arrow[ll,"\alpha"']\arrow[dl,"\alpha^*(\sigma)=\sigma\circ\alpha"]\\
\& S(X)'=\mathcal X \& {}
\end{tikzcd}}}}

% foot, end of bracketed note and final NB
et on se donne un hom de celui-ci dans le complexe analogue $A^\ast_{*\mathcal X}$ sur $\mathcal X=S(X)'$ défini par $A^\ast_{*\mathcal X}$

\uline{NB} finalement \textcircled{1} $(A_*^*,\deleted{\unsure{R|\varepsilon}}k_*,K_*)$ et (4) ne servent que pour pouvoir définir $\unsure{\hookrightarrow|\Leftarrow|\text{un}}\Omega^*\longrightarrow S^{*\,\prime}_X$ comp. \unsure{co-augm.|coaugm.|aux augm.}

% lower left
% archive sheet number: 69
